\hypertarget{lexical-analysis-and-the-execution-model}{%
\chapter{Lexical Analysis and the Execution
Model}\label{lexical-analysis-and-the-execution-model}}

Python is often described as an interpreted language, but this is a
high-level abstraction. Under the hood, CPython follows a classic
compiler-interpreter pipeline: Lexical Analysis \(\rightarrow\) Parsing
\(\rightarrow\) Abstract Syntax Tree (AST) \(\rightarrow\) Bytecode
Generation \(\rightarrow\) Virtual Machine Execution. While Chapter 1
introduced the LL(1) pipeline and Chapter 15 the PEG transition, this
chapter formalizes the execution model and the mechanics of dynamic
evaluation.

\hypertarget{lexical-analysis-from-raw-bytes-to-tokens}{%
\subsection{31.1 Lexical Analysis: From Raw Bytes to
Tokens}\label{lexical-analysis-from-raw-bytes-to-tokens}}

The first stage of execution is the \textbf{Tokenizer}. In CPython, the
tokenizer is implemented in C
(\passthrough{\lstinline!Parser/tokenizer.c!}). Its job is to break the
stream of source code (or bytes) into a stream of logical units called
\textbf{Tokens}.

\hypertarget{token-types-and-the-tokenize-module}{%
\subsubsection{\texorpdfstring{1. Token Types and the \texttt{tokenize}
Module}{1. Token Types and the tokenize Module}}\label{token-types-and-the-tokenize-module}}

Python exposes its internal tokenizer via the
\passthrough{\lstinline!tokenize!} standard library module.

\begin{lstlisting}[language=Python]
import tokenize
from io import BytesIO

code = "x = 5 + 10"
tokens = tokenize.tokenize(BytesIO(code.encode('utf-8')).readline)
for token in tokens:
    print(token)
\end{lstlisting}

Each token contains: * \textbf{Type}: \passthrough{\lstinline!NAME!},
\passthrough{\lstinline!NUMBER!}, \passthrough{\lstinline!OP!},
\passthrough{\lstinline!NEWLINE!}, \passthrough{\lstinline!INDENT!},
\passthrough{\lstinline!DEDENT!}. * \textbf{String}: The actual text
(e.g., ``x'', ``5''). * \textbf{Start/End Pos}: Line and column numbers
for error reporting.

\hypertarget{the-indentation-stack}{%
\subsubsection{2. The Indentation Stack}\label{the-indentation-stack}}

Unlike most languages, Python's tokenizer is stateful regarding
whitespace. It maintains an \textbf{Indentation Stack}. * When a line
has more leading whitespace than the top of the stack, it emits an
\passthrough{\lstinline!INDENT!} token and pushes the new level onto the
stack. * When it has less, it pops from the stack and emits one or more
\passthrough{\lstinline!DEDENT!} tokens until the levels match.

\hypertarget{the-ast-and-the-compilation-pipeline}{%
\subsection{31.2 The AST and the Compilation
Pipeline}\label{the-ast-and-the-compilation-pipeline}}

Once tokens are parsed into a tree structure (PEG parser), CPython
transforms the Concrete Syntax Tree (CST) into an \textbf{Abstract
Syntax Tree (AST)}.

\hypertarget{the-ast-module}{%
\subsubsection{\texorpdfstring{1. The \texttt{ast}
Module}{1. The ast Module}}\label{the-ast-module}}

The AST is the representation that Python's optimizer and code generator
actually use. You can inspect it using the \passthrough{\lstinline!ast!}
module:

\begin{lstlisting}[language=Python]
import ast
tree = ast.parse("x = 5 + 10")
print(ast.dump(tree, indent=4))
\end{lstlisting}

The output shows a \passthrough{\lstinline!Module!} containing an
\passthrough{\lstinline!Assign!} node, with a
\passthrough{\lstinline!Name!} target and a
\passthrough{\lstinline!BinOp!} (Add) value.

\hypertarget{constant-folding-and-peephole-optimization}{%
\subsubsection{2. Constant Folding and Peephole
Optimization}\label{constant-folding-and-peephole-optimization}}

During the transition from AST to bytecode, CPython performs simple
optimizations: * \textbf{Constant Folding}: Expressions like
\passthrough{\lstinline!1 + 2!} are evaluated at compile-time and
replaced with \passthrough{\lstinline!3!}. * \textbf{Dead Code
Elimination}: Code following a \passthrough{\lstinline!return!} or
\passthrough{\lstinline!raise!} that is unreachable is stripped.

\hypertarget{dynamic-execution-eval-vs.-exec}{%
\subsection{\texorpdfstring{31.3 Dynamic Execution: \texttt{eval()}
vs.~\texttt{exec()}}{31.3 Dynamic Execution: eval() vs.~exec()}}\label{dynamic-execution-eval-vs.-exec}}

Python provides two primary built-ins for dynamic code execution. Their
difference lies in what they accept and what they return.

\hypertarget{evalexpression-globalsnone-localsnone}{%
\subsubsection{\texorpdfstring{1.
\texttt{eval(expression,\ globals=None,\ locals=None)}}{1. eval(expression, globals=None, locals=None)}}\label{evalexpression-globalsnone-localsnone}}

\begin{itemize}
\tightlist
\item
  \textbf{Input}: A single Python expression (something that can be on
  the right side of an assignment).
\item
  \textbf{Output}: The value of the expression.
\item
  \textbf{Internals}: Compiles the string to a code object with the
  \passthrough{\lstinline!eval!} mode, then executes it in the provided
  namespaces.
\end{itemize}

\hypertarget{execobject-globalsnone-localsnone}{%
\subsubsection{\texorpdfstring{2.
\texttt{exec(object,\ globals=None,\ locals=None)}}{2. exec(object, globals=None, locals=None)}}\label{execobject-globalsnone-localsnone}}

\begin{itemize}
\tightlist
\item
  \textbf{Input}: A block of Python code (statements, class/function
  definitions).
\item
  \textbf{Output}: Always \passthrough{\lstinline!None!}.
\item
  \textbf{Internals}: Compiles the code in
  \passthrough{\lstinline!exec!} mode. It modifies the
  \passthrough{\lstinline!locals!} dictionary (if provided) to include
  newly defined variables.
\end{itemize}

\hypertarget{the-compile-function-and-code-objects}{%
\subsection{\texorpdfstring{31.4 The \texttt{compile()} Function and
Code
Objects}{31.4 The compile() Function and Code Objects}}\label{the-compile-function-and-code-objects}}

Both \passthrough{\lstinline!eval!} and \passthrough{\lstinline!exec!}
use \passthrough{\lstinline!compile()!} under the hood. For performance,
you should pre-compile code if you intend to run it multiple times.

\begin{lstlisting}[language=Python]
code_str = "print('Hello, Godhood')"
# Modes: 'exec' for blocks, 'eval' for expressions, 'single' for REPL-style
code_obj = compile(code_str, filename="<string>", mode="exec")

# Inspection
print(code_obj.co_code)      # Raw bytecode
print(code_obj.co_consts)    # Constants used in the code
print(code_obj.co_names)     # Global/Builtin names used
\end{lstlisting}

\passthrough{\lstinline!PyCodeObject!} is the C struct that holds this
data. When \passthrough{\lstinline!exec(code\_obj)!} is called, the
CPython VM pushes a new \passthrough{\lstinline!PyFrameObject!} onto the
evaluation stack and hands the code object to the interpreter loop
(\passthrough{\lstinline!\_PyEval\_EvalFrameDefault!}).

\begin{center}\rule{0.5\linewidth}{0.5pt}\end{center}

\begin{center}\rule{0.5\linewidth}{0.5pt}\end{center}
